% Transcription — fonds Alexandre Grothendieck, shelfmark 6, batch 2 (pages 21-40).
% Established from the facsimile archives/batches/6/batch-02.pdf, cut from a
% copy of 6.pdf fetched through the project's own relay
% (https://grothendieck-relay.onrender.com/source/6.pdf) rather than from
% grothendieck.umontpellier.fr directly; per the skill transcribe-grothendieck.
% The batch file carries no cover sheet: PDF page k is the archivists' page
% 20+k. Their pencil numbers bottom-left agree from 21 to 38; the two last
% leaves are pencilled 50 and 51 although the facsimile places them at 39 and
% 40 (the folder was evidently re-ordered after pencilling); \page{} follows
% the facsimile's order, and a \note says so.
%
% Pass: Opus 5.5 (claude-opus-5-5), 2026-09-26 — first pass, unchecked against the pages by a human.
%
% The folder sits in the inventory group "4-9" « Cristaux »; it has its own
% date in the catalogue, which \dating{} carries verbatim, with no group
% sentence.
%
% What the batch contains, and the decisions taken.
% - Pages 21-31: the end of HIS OWN corrected typed copy of the May 1966
%   letter to Tate, continuing batch 1 mid-word (« condi-/tion », §2.3). Same
%   conventions as batch 1: typed text in full; his ink (and pencil)
%   interventions as \add / \struck with French notes; machine strike-overs
%   as \struck (\ill when unreadable); typescript slips kept with \note{sic};
%   his letter page number in a \note at each page (letter p. = archivists'
%   p. - 1: pp. 20-30 of the letter). The letter ends on p. 31 (letter
%   p. 30) in the middle of §3.2, Chapter 3 « Remarques sur les groupes
%   p-divisibles », with no closing formula or signature; the rest of that
%   page is blank.
% - Pages 32-36: HIS OWN MANUSCRIPT NOTES, in ink on white paper, under a
%   cover sheet (p. 32) inscribed in his hand in pencil « Réduction des
%   courbes elliptiques et T_p(E) (cf. aussi livre de Serre "Abelian l-adic
%   representations") ». Transcribed, as his mathematics; the cover gets no
%   \page, its words become a section heading with a \note. These are
%   fast drafts: formulas come through, much connective prose is \ill.
% - Page 37: a cover in his hand, « Gribouillis de Tate sur sa théorie »:
%   heading only, no \page.
% - Page 38: notes in ANOTHER HAND (by the cover, Tate's), in English, very
%   faded ink (L/K Galois, H^1(G, R_L), Thm.). Someone else's document with
%   no mark of his: identified in one French \note under \page{38}, not
%   re-set.
% - Page 39 (pencilled 50): a cover in his hand, « Questions diverses en
%   théorie de Tate-Hodge »: heading only, no \page.
% - Page 40 (pencilled 51): his ink list of questions a)-f) on Tate-Hodge
%   theory, written across an unrelated typed page (show-through of a French
%   typescript, not reproduced). Transcribed.
%
% Pages skipped (no \page): 32, 37, 39 — covers, whose words are carried as
% section headings. Page noted only: 38 (not his).
%
% What the letter pages are about. §2.3 end: the frobenius F_p on DR(f)_p
% via the relative frobenius, and V by duality with FV = VF = p^{2d}, so the
% De Rham isocrystal is a bi-isocrystal; the admission that DR^i is not a
% bicrystal of weight i in general. §2.4: abelian schemes: DR^1(A/S) a
% Dieudonné crystal; the three « affirmations » (DR^1 depends only on A and
% S -> T; over perfect k it is the classical Dieudonné module; lifts of A
% correspond to lifts of the Hodge filtration). §2.5-2.6: Serre's universal
% vector extension G(A) of A by V(t_{A*}), its Lie algebra as the dual of the
% Hodge sequence, the crystal in groups G(A/S); NB on Manin's map and
% Picard-Fuchs. What he checked: char. 0, rigidity of G(A), first-order
% lifts. §2.7: « je m'aperçois avec consternation » — counterexample from
% supersingular (Hasse-invariant zero) elliptic curves; one must pass to
% isocrystals. §2.8: ordinary abelian varieties, Serre's canonical lifts, the
% modular curve over Z_p; the notion of crystal « irrémédiablement trop
% fine » for families. Ch. 3 (3.1-3.2): the same construction for a
% p-divisible group (written with the typed sign Ø, set as \varnothing):
% universal extension G(Ø), its Lie algebra H(Ø) in 0 -> t^v_{Ø*} -> H(Ø)
% -> t_Ø -> 0, and lifts over a complete DVR as lifts of a filtration of
% M ⊗_W A. The manuscript notes (pp. 33-36): p-adic Tate modules of
% elliptic curves over a local field — ordinary good reduction (extension of
% L by L^v(1), class in H^1(K, L^{-2} ⊗ T_p(G_m)) ≃ Hom(L_0^2, K^*^)),
% supersingular reduction, potentially good and multiplicative reduction
% (Tate-curve extension of Z_p by Z_p(1)); p. 40: questions on the Hodge-Tate
% comparison H_DR(X_C) ≃ H_Hdg(X_C), Tate's conjecture, Washnitzer-Monsky.
%
% Notation fixed once for the batch, as in batch 1: \underline for letters
% underlined at the machine (Z, O, Ω, R, G, M, t, H); DR, Topcr roman; the
% typed Ø as \varnothing; ~> as \rightsquigarrow. In the manuscript notes:
% bold G_m, Z_p, Q_ℓ; \check{L} and L^{-2} as he writes them.
%
% Readings to check first: the whole of pp. 33-36 (fast hand; e.g.
% « Einseinheiten » p. 34, « tordue » p. 36, the bracketed pencil note top
% right of p. 33); the pencil marginalia on pp. 26-27 (« faux », « à
% éclaircir … », « faux, c'est jamais vrai en car p>0 ! »); the struck
% words of the typescript left \ill.

\documentclass[11pt,a4paper]{article}
% Le préambule commun à toutes les transcriptions du fonds Grothendieck.
%
% Il tient en une idée : **l'incertitude se compose, elle ne se résout pas.**
% Une transcription qui lisse un mot douteux a détruit la seule chose pour
% laquelle on la faisait — un lecteur doit pouvoir savoir, à chaque endroit,
% ce qui était sur le feuillet et ce qui a été ajouté. Les six macros qui
% suivent sont donc l'appareil critique, et non de la décoration.
%
% Les mêmes commandes sont reconnues par `scripts/render.mjs`, qui produit la
% vue de lecture du site. Une macro ajoutée ici doit l'être aussi là-bas,
% faute de quoi le rendu échouera bruyamment — ce qui est voulu : un rendu
% silencieusement faux est pire qu'un rendu absent.

\ProvidesPackage{grothendieck}[2026/08/08 Transcriptions du fonds Grothendieck]

\usepackage{amsmath,amssymb,amsthm}
\usepackage{mathrsfs}
\usepackage{stmaryrd}
% Les diagrammes commutatifs sont partout dans ce fonds : c'est la langue
% même de la géométrie algébrique de Grothendieck, et une transcription qui
% les rendrait en prose ne transcrirait rien.
\usepackage{tikz-cd}
% Français par défaut — c'est la langue des feuillets — mais l'anglais est
% chargé aussi : la traduction et le résumé sont des documents anglais, et la
% typographie française (espaces avant les deux-points) y serait une faute.
% Ces documents basculent par \selectlanguage{english} en tête de corps.
\usepackage[english,french]{babel}
\usepackage{fontspec}
\usepackage{geometry}
\geometry{a4paper,margin=25mm,marginparwidth=28mm}
\usepackage{marginnote}
\usepackage{xcolor}
% ulem, not soul. `soul`'s \st and \ul are built on a character-by-character
% scan that XeLaTeX's font machinery breaks: the document fails with an
% undefined control sequence rather than degrading. ulem does the same two jobs
% and is Unicode-engine safe. `normalem` stops it hijacking \emph, which the
% transcriptions use for Grothendieck's own emphasis.
\usepackage[normalem]{ulem}
\usepackage{hyperref}
\hypersetup{colorlinks=true,linkcolor=black,urlcolor=black,pdfborder={0 0 0}}

% --- Métadonnées du lot -----------------------------------------------------
% Reprises telles quelles par le script de rendu, qui les affiche en tête de
% la vue de lecture. Elles fixent ce que le fichier prétend couvrir : sans
% elles, une transcription flotte sans point d'attache dans un fonds de seize
% mille feuillets.

\newcommand{\folder}[1]{\gdef\grothFolder{#1}}
\newcommand{\batch}[1]{\gdef\grothBatch{#1}}
\newcommand{\leaves}[2]{\gdef\grothFirst{#1}\gdef\grothLast{#2}}
\newcommand{\foldertitle}[1]{\gdef\grothFoldertitle{#1}}
\gdef\grothFolder{?}\gdef\grothBatch{?}\gdef\grothFirst{?}\gdef\grothLast{?}\gdef\grothFoldertitle{}

% --- Le résumé --------------------------------------------------------------
%
% Il ouvre la lecture modernisée, sous le titre « Résumé », et il est composé
% en petit. Ce n'est pas un ornement : c'est le seul passage du document qui
% ne soit pas des mathématiques, et un lecteur doit pouvoir le reconnaître
% avant de l'avoir lu — puis le sauter s'il connaît déjà le sujet, ou s'y
% arrêter s'il ne le connaît pas. Le filet qui le suit marque la reprise du
% corps.

\newenvironment{resume}
  {\small\setlength{\parskip}{0.45em}}
  {\par\medskip\hrule\medskip}

% --- Mots-clés ---------------------------------------------------------------
%
% En anglais, en clôture du résumé de la lecture modernisée : le vocabulaire
% moderne sous lequel chercher ce que les feuillets construisent. Le manifeste
% du site (`npm run manifest`) les extrait pour étiqueter la cote entière —
% cette ligne est la source unique des tags, il n'y a pas de fichier à côté.
\newcommand{\keywords}[1]{\par\smallskip\noindent
  {\footnotesize\textsc{Keywords} --- \textit{#1}}\par}

% --- L'appareil critique ----------------------------------------------------

% Le numéro de feuillet, en marge. C'est l'ancre qui permet de régler un
% désaccord sur une formule en regardant *un* feuillet plutôt qu'une chemise
% de six cents. Le site s'en sert aussi pour tourner les pages du fac-similé
% quand on fait défiler la transcription.
\newcommand{\leaf}[1]{%
  \marginnote{\footnotesize\color{gray!70}\textsf{#1}}%
  \label{leaf:#1}\ignorespaces}

% Illisible : on ne devine pas. Un mot inventé qui se lit comme les autres est
% le pire résultat possible.
\newcommand{\ill}{\textcolor{red!60!black}{[\,\ldots\,]}}

% Lecture douteuse : le mot est donné, et signalé comme incertain.
\newcommand{\uncertain}[1]{\uline{#1}}

% Ajout de l'éditeur — une lettre manquante, un mot sous-entendu.
\newcommand{\add}[1]{\textcolor{blue!50!black}{[#1]}}

% Biffé par Grothendieck lui-même. Ce qu'il a rayé fait partie du document :
% une rature dit souvent où la pensée a changé de direction.
\newcommand{\struck}[1]{\sout{#1}}

% Note du transcripteur, sur l'état matériel du feuillet.
\newcommand{\note}[1]{\par\smallskip{\small\color{gray!60!black}\textit{[#1]}}\par\smallskip}

% Note marginale de Grothendieck, distincte du corps du texte.
\newcommand{\marginal}[1]{\par\smallskip\noindent\hspace{1em}%
  {\small\color{gray!60!black}#1}\par\smallskip}

% --- Titre ------------------------------------------------------------------

\AtBeginDocument{%
  \begin{center}
    {\large\bfseries Fonds Alexandre Grothendieck --- cote n\textsuperscript{o}~\grothFolder}\\[2pt]
    {\itshape\grothFoldertitle}\\[4pt]
    {\small feuillets \grothFirst--\grothLast\ (lot \grothBatch)}\\[2pt]
    {\footnotesize\color{gray!60!black}Transcription établie d'après le fac-similé numérisé
     par l'Université de Montpellier.\\
     Les lectures incertaines sont soulignées, les illisibles notées [\,\ldots\,],
     les ajouts entre crochets.}
  \end{center}
  \vspace{1em}}

\endinput


\folder{6}
\batch{2}
\pages{21}{40}
\dating{1966-[à partir de 1970]}
\watermark{Édition de démonstration}
\foldertitle{Lettre Tate (mai 66) (Cristaux) : lettre (1966), tapuscrit (s.d.), notes manuscrites (s.d.)}

\begin{document}

\page{21}
\note{p.~20 de la lettre. Le mot « condi-/tion » est coupé entre les deux
lots.}
tion de compatibilité avec le changement de base, on constate que
$\mathrm{DR}(f_p)^{(p)}$ \struck{\ill{}} n'est autre que
$\mathrm{DR}(f_p^{(p)})$, où $f_p^{(p)}\colon X_p^{(p/S)}\to S_p$ est le
morphisme structural du frobeniusé relatif de $X_p$ sur $S_p$. Or on a le
morphisme de frobenius $X_p^{(p/S)}\to X$, qui est un
$\underline{S}$-morphisme, et par fonctorialité de $\mathrm{DR}$ nous donne
\[
  \mathrm{DR}(f)_p^{(p)}\longrightarrow\mathrm{DR}(f)_p
\]
comme on voulait.\note{l'indice $p$ de « $S_p$ » (après la flèche) est
récrit à l'encre ; le « $\mathrm{DR}(f)_p^{(p)}\to\mathrm{DR}(f)_p$ » est dans
le cours du texte tapé. La flèche de frobenius $X_p^{(p/S)}\to X$ : tel quel
(on attendrait $X_p$ au but).} Il faut prouver que cet homomorphisme est en
fait une isogénie, donc que l'\struck{bi}isocristal défini par
$\mathrm{DR}(f)$ devient, à l'aide des $F_p$, un bi-isocristal.\note{« bi »
est barré d'un trait oblique à l'encre.} Mais utilisant la relation de
dualité écrite dans 2.2.\ (à vrai dire, l'écriture $T^d$ pour le cristal unité
ne prend son sens que lorsque on le regarde comme muni de sa structure de
\emph{bi}-cristal naturelle, qui n'intervient qu'ici), on peut transposer $F$
en un $V$, tel que $FV=VF=p^{2d}$, ce qui prouve notre assertion.\note{« bi »
souligné deux fois à l'encre.}

D'ailleurs, lorsque l'on passe du cristal de De Rham $\mathrm{DR}(f)$ à
l'isocristal correspondant, \struck{et donc} des objets de cohomologie
$\mathrm{DR}^i(f)$ aux isocristaux correspondants, \struck{\ill{}} il
\struck{\ill{}} \add{sera} vrai (tout comme pour les
$\underline{R}^if_*$ de De Rham en caractéristique nulle) que leur formation
commute à tout changement de base, de sorte que chacun des isocristaux
$\mathrm{DR}^i(f)$ devient à son propre titre un bi-isocris-
tal. Au moment de rédiger 1.10.\ il m'avait semblé\add{ que}, sans mettre du
iso dans le coup, \struck{que} $\mathrm{DR}^i(f)$ devrait être \add{un}
bi-cristal de poids $i$, mais je m'étais canulé, il faudrait pour cela une
polarisation de $X$ sur $S$ qui définisse un isomorphisme (pas seulement une
isogénie) de $\mathrm{DR}^i(f)$ avec $\mathrm{DR}^{2d-i}\otimes T^{-(d-i)}$,
ce qui n'existe évidemment que très exceptionellement ; une fois qu'on l'a, on
définit $V$ dans $\mathrm{DR}^i$ en transposant $F$ dans
$\mathrm{DR}^{2d-i}$.\note{« que » et « un » sont à l'encre au-dessus de la
ligne, le « que » tapé devant $\mathrm{DR}^i(f)$ est biffé à l'encre.}

2.4. \emph{Cas d'un schéma abélien.} Il semble que ce ne soit guère que dans ce
cas-là où il soit bien raisonnable de parler de \add{bi-}cristaux et non
seulement de bi-isocristaux. De façon précise, $\mathrm{DR}^1(f)$ a dans ce
cas une structure de bicristal, défini par les $F_p$ précédents\add{,} et des
$V_p$ qui se définissent encore, par fonctorialité de $\mathrm{DR}$, à l'aide de
l'homomorphisme « Verschiebung » $A_p^{(p/S_p)}\to A_p$ (défini
avec la généralité qui convient dans les notes de Gabriel dans SGAD). Comme

\page{22}
\note{p.~21 de la lettre.}
$\mathrm{DR}^1(A/S)$ est un foncteur \struck{\ill{}} multiplicatif en $A$,
grâce à Kunneth\note{sic.} postulé dans 2.2., et que l'on a
$FV=VF=p.\,\mathrm{id}$ sur les schémas abéliens en car.\ $p$, on en conclut
les mêmes relations dans $\mathrm{DR}^1$. Cela montre donc que
$\mathrm{DR}^1(A)$ est un bicristal de poids 1, i.e.\ un \emph{cristal de
Dieudonné}. Lorsque $S$ est le spectre d'un corps \add{parfait} de
caractéristique $p>0$, un tel cristal s'identifie simplement, d'après 1.12, à
un module de Dieudonné sur $W=W(k)$. Bien sûr, \emph{on doit trouver que ce
module n'est autre que le module de Dieudonné du groupe formel} \struck{\ill{}}
(ou plutôt, du groupe $p$-divisible) \emph{défini par la variété abélienne}
$A$.\note{« parfait » est tapé dans l'interligne et appelé à l'encre ; le « c »
de « bicristal » est corrigé à l'encre.} Débarassé\note{sic.} de
l'hyperstructure axiomatico-conjecturale de 2.2., \struck{en fait} on peut
exprimer ainsi la construction \struck{ainsi} obtenue du module de Dieudonné :

1º) Soit $A$ un schéma abélien sur un schéma affine $S$. On sait que pour tout
morphisme $S\to T$ d'immersion fermée surjective, défini par Idéal nilpotent
sur $T$, $A$ se prolonge en un schéma abélien $B$ sur $T$. On peut regarder la
cohomologie de De-Rham ordinaire
$\underline{R}^1g_*(\underline{\Omega}^{\bullet}_{B/T})$, où $g\colon B\to T$
est le morphisme structural, et on sait que c'est un Module localement libre de
rang $2g$, où $g$ est la dimension relative de $A$ sur $S$ donc de $B$ sur $T$.
\emph{Première affirmation} : ce Module, à isomorphisme canonique près, ne
dépend que \add{de $A$ sur $S$ et} de $S\to T$, et pas du choix du prolon-
gement $B$ de $A$. En d'autres termes, la donnée de $A$ sur $S$ définit un
\emph{cristal} localement libre de rang $2g$ sur $S$, soit
$\mathrm{DR}^1(A/S)$. C'est évidemment un foncteur additif en $A$, compatible
avec le changement de base\add{, et définissable sans condition affine sur $S$
(en recollant).}\note{« de $A$ sur $S$ et » (fin de la page précédente) et
« , et définissable \ldots\ (en recollant). » sont à l'encre, dans
l'interligne ; le « t » de « définit » aussi.}

2º) Supposons que $S$ soit le spectre d'un corps parfait $k$ de car $p>0$,
alors le cristal précédent s'identifie à un module libre de rang $2g$ sur
$W=W(k)$. On y introduit les structures $F$ et $V$, en utilisant comme
ci-dessus les homomorphismes $F\colon A\to A^{(p)}$ et $V\colon A^{(p)}\to A$.
On trouve ainsi un module de Dieudonné \add{$M$}, et : \emph{Deuxième
affirmation} : C'est bien celui défini par Dieudonné. En d'autres termes : si
$\Lambda$ est n'importe quel anneau local artinien de corps résiduel $k$, et
$B$ un prolongement de $A$ en un schéma abélien sur $\Lambda$, alors
\struck{\ill{}} la cohomologie de De Rham
\[
  H^1_{\mathrm{DR}}(B)=\underline{H}^1(B,\underline{\Omega}^{\bullet}_{B/\Lambda})
\]
est canoniquement \add{et fonctoriellement} isomorphe à $M\otimes_W\Lambda$, où
$M$ est le module de Dieudonné classique de $A$, et où on tient compte du th.\
de Cohen qui munit $\Lambda$ d'une structure canonique de
$W$-algèbre.\note{« $M$ » est à l'encre ; « et fonctoriellement » est tapé
dans l'interligne et appelé à l'encre. L'anneau artinien est tapé « A » puis
corrigé en $\Lambda$.} (NB la \struck{canonicité} fonctorialité de

\page{23}
\note{p.~22 de la lettre.}
l'isomorphisme garantira automatiquement qu'il est compatible avec $F$ et $V$).

\add{2.5.} Je n'ai pas vérifié, à vrai dire, les deux affirmations, mais n'ai
pas le moindre doute qu'elles sont correctes telles quelles. Cette façon de
voir le module de Dieudonné permet de plus d'expliciter de façon remarquable
les variations infinitésimales de structure de la variété abélienne $A$ donnée
en caractéristique \struck{nulle} $p>0$, en termes du module de Dieudonné :
les prolongements de $A$ en un schéma abélien $B$ sur $\Lambda$ doivent
correspondre exactement aux \struck{\ill{}} modules quotients libres de rang
$g$ de $M\otimes_W\Lambda$ qui redonnent, modulo $p$, le module quotient de
rang $g$ canonique de $M\otimes_W k=H^1_{\mathrm{DR}}(A)$, \add{(}correspondant
à la filtration canonique de cette cohomologie\add{)}. Plus généralement et
plus précisément :\note{le numéro « 2.5. » et les parenthèses sont à
l'encre.}

3º) Considérons, pour tout schéma abélien $A$ sur une base $S$, sur la
cohomologie de De Rham
$R^1f_*(\underline{\Omega}^{\bullet}_{A/S})=H^1(f)$, la filtration canonique
\[
  0\longleftarrow R^1f_*(\underline{O}_A)\longleftarrow H^1(f)\longleftarrow
  R^0f_*(\underline{\Omega}^1_{A/S})\longleftarrow 0\ .
\]
Donc pour tout $B$ sur $T$ comme dans 1º, on a sur
\add{$\underline{M}(T)=$} $\mathrm{DR}^1(A/S)(T)=H^1(g)$ une filtration
naturelle, ne dépendant que de la classe à isomorphisme près du relèvement
choisi $B$ de $A$ à $T$, et prolongeant la filtration déjà connue de
\struck{$\mathrm{DR}^1$} $\underline{M}(S)=H^1(f)$\note{« $\underline{M}(T)=$ »
est à l'encre dans l'interligne ; « $\underline{M}(S)$ » est récrit à l'encre
sur un mot tapé biffé.}
provenant de $A$. Ceci dit, \emph{troisième affirmation} : \emph{on obtient
ainsi une correspondance biunivoque entre classes de prolongements de} $A$
\emph{à} $T$, \emph{et prolongements de la filtration donnée de}
$\underline{M}(S)=\underline{M}(T)\otimes_{\underline{O}_T}\underline{O}_S$
\emph{en une filtration de} $\underline{M}(T)$. Plus précisément encore, le
foncteur $B\rightsquigarrow(A,\varnothing)$, de la catégorie des schémas
abéliens \add{$B$} sur $T$, dans la catégorie des schémas abéliens
\struck{\ill{}} $A$ sur $S$, munis d'une filtration \add{$\varnothing$} de
$\mathrm{DR}^1(A/S)(T)$ prolongeant celle de
$\mathrm{DR}^1(A/A)(S)$\note{sic, pour $\mathrm{DR}^1(A/S)(S)$.} (NB il ne
s'agit que de filtrations à quotients localement libres, bien sûr) est une
équivalence de catégories.\note{« $B$ » et « $\varnothing$ » (la filtration)
sont à l'encre, dans l'interligne.}

Cet énoncé\note{tapé « Ces énoncés » et corrigé à la machine.} est
évidemment fort suggestif aussi pour des généralisations en cohomologie de De
Rham de dimension supérieure, pour un morphisme lisse et projectif quelconque,
tenant compte de la filtration canonique de celle ci. On voit bien en tous cas
que ce dernier élément de structure n'est \emph{pas} de nature cristalline,
i.e.\ donnée par une filtration du cristal de De Rham postulé dans 2.2., mais
est au contraire dans la nature d'un élément de structure « conti-

\page{24}
\note{p.~23 de la lettre.}
nu », dont la variation doit refléter fidèlement les variations de structure du
motif donnant naissance au cristal envisagé. Pour arriver à préciser ce dernier
point, il faudrait des fondements un peu plus fermes de la théorie des motifs,
comme de celle (certainement beaucoup plus élémentaire) des cristaux et de la
cohomologie de De Rham. Une autre généralisation, (suggérée par comparaison du
3º avec Serre-Tate, disant que si $S$ est de caractéristique $p$ \add{$>0$},
alors étendre $A$ à $T$\add{,} c'est pareil qu'étendre le groupe $p$-divisible
correspondant), concerne la théorie des groupes formels, dont il sera question
au Chapitre 3. Pour préparer le terrain, je vais présenter d'une façon un peu
différente le cristal $\mathrm{DR}^1(A/S)$ associé à un schéma abélien, en
utilisant explicitement la structure de groupe dudit.\note{« $>0$ » et la
virgule après « $T$ » sont à l'encre.}

2.6. De façon générale, paraphrasant sur une base quelconque une vieille
construction de Serre (c'est de lui que je l'ai apprise, du moins), on trouve
que pour tout schéma abélien $A$ sur une base $S$, il y a une extension
naturelle de $A$ par le fibré vectoriel
$V(R^1f_*(\underline{O}_A))=V(\underline{t}_{A}\ )$, (où $A$ est le schéma
abélien dual de $A$).\note{un blanc est laissé après l'indice de
$\underline{t}_A$ et après le second « $A$ » : le signe du dual, à ajouter à la
main, n'a pas été ajouté ; on attend $\underline{t}_{A'}$, « où $A'$ est le
schéma abélien dual ».} Attention à la notation, le fibré vectoriel
$V(\underline{E})$ est contravariant en $\underline{E}$, ces sections sont les
homomorphismes $\underline{E}\to\underline{O}_S$ ! L'extension en question est
universelle parmi les extensions de $A$ par des fibrés vectoriels, est
fonctorielle en $A$, et compatible avec changement de base. Appelons la
$G(A)$. Ainsi, le faisceau de Lie de $G(A)$ est une extension
\[
  \begin{array}{ccccccc}
  0\to & R^1f_*(\underline{O}_A)^{\vee} & \to & \underline{g}(A) & \to &
  R^0f_*(\underline{\Omega}^1_{A/S})^{\vee} & \to 0\\
  & \| & & & & \| & \\
  & \underline{t}^{\vee}_{A^*} & & & & \underline{t}_A &
  \end{array}\ ,
\]
\note{l'astérisque de $A^*$ est à l'encre.}
qui est duale de la suite exacte envisagée dans 2.5., dont nous avons donc ici
une construction indépendante en termes d'extensions de $A$ par des groupes
vectoriels. On peut en profiter pour préciser encore les assertions 1º et 3º
de 2.4.\ et 2.5. L'assertion 1º se précise en affirmant que, pour un
prolongement $B$ de $A$, de $S$ à $T$, le schéma en groupes $G(B/S)$ est
déterminé, à isomorphisme canonique \add{(et même \emph{unique})} près, par la
seule donnée de $A$ et de \struck{l'extension} $S\to T$, indépendamment du
choix particulier du prolongement $B$. En d'autres termes, on trouve un
\emph{cristal}\note{« (et même unique) » est à l'encre dans l'interligne,
« unique » souligné ; « l'extension » est biffé à la machine. $G(B/S)$ :
tel quel (on attendrait $G(B/T)$).}

\page{25}
\note{p.~24 de la lettre.}
\emph{en schémas de groupes} \add{lisses} $\underline{G}(A/S)$, et pour chaque
prolongement \add{infinitésimal} $B$ de $A$ à un $T$, un isomorphisme canonique
$\underline{G}(A/S)(T)\simeq G(B/T)$ ; tout ça bien sûr fonctoriel en $A$ et
compatible avec changements de base. D'autre part, on peut préciser alors 3º
en indiquant quel est le foncteur quasi-inverse de celui envisagé dans cet
énoncé : l'extension de la filtration de $\mathrm{DR}^1(A/S)(S)$ en une
filtration du faisceau de Lie $\mathrm{DR}^1(A/S)(T)^{\vee}$ de
$\underline{G}(A/S)(T)$ revient à étendre le sous-groupe vectoriel canonique de
$\underline{G}(A/S)(S)$ en un sous-groupe vectoriel de
$\underline{G}(A/S)(T)$, et l'on trouve $B$ en passant simplement au
quotient.\note{« lisses » et « infinitésimal » sont tapés dans l'interligne,
le second appelé à l'encre ; le $\vee$ de $\mathrm{DR}^1(A/S)(T)^{\vee}$ est à
l'encre.}

NB Je suis tombé sur la connexion canonique de $G(A/S)$ en essayant de
simplifier la construction de Manin de l'application
$A(S)\to\Gamma(S,\underline{O}_S)$ associée à une équation de
Picart\note{sic.}-Fuchs sur $S$\add{,} relativement à $A$. Pour ceci, il suffit
de noter que localement sur $S$ toute section de $A$ sur $S$ se remonte en une
section de $G=G(A/S)$ sur $S$. La donnée de la connexion de $G$ permet alors
de prendre la dérivée de cette section, qui est un élément de
$\Gamma(S,\underline{t}_G\otimes\underline{\Omega}^1_S)$, déterminé modulo
une section de l'image du faisceau $\underline{t}^{\vee}_{A^*}$, correspondant
à l'indétermination du relèvement d'\emph{une} section de $A$ en une section de
$G$. Les équations de Picard-Fuchs sont définies tout juste pour arriver, d'une
telle section de\note{le $\vee$ et l'astérisque de
$\underline{t}^{\vee}_{A^*}$ sont à l'encre ; le « d » de « Picard » est
corrigé à l'encre sur un « t » tapé ; « une » est souligné à l'encre.}
\struck{\ill{}} $\underline{t}_G\otimes\underline{\Omega}^1_S$ (qui en fait est
un \add{co}\emph{cycle}, compte tenu de la connexion canonique de
$\underline{t}_G$ provenant de $G$), avec l'indétermination qu'on vient de
préciser, à tirer une section de $\underline{O}_S$ \ldots\ (Du moins en
caractéristique nulle, cas dans lequel se place Manin de toutes
façons).\note{« co » est à l'encre devant « cycle », souligné à la machine.
Le reste de la page est blanc.}

\page{26}
\note{p.~25 de la lettre.}
Voici les résultats positifs que j'ai vérifiés dans la direction des
assertions précédentes :

a) Si $S$ est de caractéristique nulle, les assertions 1º et 3º (sous la forme
précisée par l'introduction du schéma en groupes $G(A)$) sont vraies.

b) Sans restrictions sur $S$, il est vrai que $G(A)$ n'a pas d'automorphismes
infinitésimaux, donc si pour deux relèvements infinitésimaux \add{donnés}
$B,B'$ de $A$, $G(B)$ et $G(B')$ sont isomorphes, l'isomorphisme entre eux est
unique.\marginal{\uncertain{faux}}\note{au crayon, dans la marge de gauche, à
hauteur de « infinitésimaux, donc si » ; lecture douteuse.} Donc si
l'\struck{assertion} \add{\uncertain{hyp.}} précédente est \struck{vrai}
\add{vérifiée}, quels que soient les relèvements infinitésimaux de $A$ au
dessus d'un ouvert $U$ de la base $S$, alors les $G(B)$ définissent
effectivement un cristal \add{en groupes} sur $S$, à fortiori les $H(B)$
définissent un cristal de modules, et $G(A)$ et $H^1_{\mathrm{DR}}(A/S)$ sont
munis de stratifications absolues canoniques, fonctorielles d'ailleurs en $A$
satisfaisant aux conditions envisagées, et compatible avec tout changement de
base qui invarie notre hypothèse sur $A$.\note{« donnés » et « en groupes »
sont tapés dans l'interligne et appelés à l'encre ; « vérifiée » est tapé
au-dessus de « vrai » biffé à l'encre ; le mot au-dessus de « assertion »,
biffé à l'encre, est de sa main.}

c) Les assertions 1º à 3º (sous la forme précisée par $G(A)$) sont vraies pour
les relèvements infinitésimaux \emph{d'ordre} 1, ou tout au moins lorsque $T$
est de la forme $D(\underline{J})$, schémas des nombres duaux défini par un
$\underline{J}$ quasi-cohérent.\marginal{à éclaircir, \ill{} peut être
\uncertain{faux}}\note{au crayon, dans la marge de gauche, en biais, à
hauteur de l'alinéa c), qu'un crochet au crayon embrasse.}

2.7. Arrivé à ce point de mes brillantes conjectures, je m'aperçois avec
consternation qu'elles sont fausses telles quelles, malgré les indications
concordantes militant en leur faveur. De façon précise, soit $k$ un corps de
caractéristique \struck{nulle} $p>0$. Je dis qu'\emph{il n'est pas possible},
\emph{pour tout} $k$-\emph{schéma} $S$ \emph{de type fini}, \emph{de dimension}
$\leqslant 1$, \emph{avec} $S_{\text{rég}}$ \emph{régulier, et tout schéma
elliptique} $A$ \emph{sur} $S$, \emph{de mettre sur} $H^1_{\mathrm{DR}}(A/S)$
\emph{une stratification relativement à} $k$, \emph{qui soit fonctorielle en}
$A$ \emph{et compatible avec les changements de base}.\marginal{\ill{}
$\leqslant 0$ !}\note{à l'encre, entouré, dans la marge de gauche à hauteur
de « de dimension $\leqslant 1$ » ; le « 1 » est barré d'un trait oblique à
l'encre, prolongé jusqu'au $A$ de $H^1_{\mathrm{DR}}(A/S)$ : sans doute pour
remplacer la dimension $\leqslant 1$ par $\leqslant 0$. Le premier mot de la
note marginale n'est pas lu.} L'ennui, comme d'habitude, provient des courbes
elliptiques de Hasse nul. Appliquant en effet les deux fonctorialités
postulées, on trouve que l'homomorphisme « frobenius » qui va de $M^{(p)}$
dans $M$ ($M=H^1_{\mathrm{DR}}(A/S)$) serait compatible avec la
stratification. Si alors $s\in S$ est un point correspondant à un $A_s$ de
Hasse nul, i.e.\ tel que frobenius sur $H^1_{\mathrm{DR}}(A_s)$ soit de carré
nul, on en conclurait aisément, grimpant sur les voisinages infinitésimaux de
$s$, que la même propriété serait vraie aux points voisins de $s$, ce qui est
évidemment faux p.ex.\ pour la famille modulaire.

\page{27}
\note{p.~26 de la lettre.}
Ceci montre que décidément, il faut en rabattre, et que dans le titre du
Chapitre et les considérations de 2.2., c'est à condition de prendre partout
des isocristaux qu'il reste une chance d'une théorie du genre de celle
envisagée précédemment. Il est fort possible d'ailleurs que la notion de
isocristal que j'ai adoptée est encore trop restrictive, en ce sens que dans
la description de 1.13.\ il faudra peut-être prendre des stratifications en
caractéristique nulle qui ne se prolongeraient pas nécessairement sur le
schéma formel sur $W$ = vecteurs de Witt tout entier. C'est une analyse
soigneuse du calcul-clef de Washnitzer-Monsky qui devrait permettre de tirer
cette question au clair.\note{un double trait vertical au crayon longe dans la
marge de gauche tout ce passage, de « Chapitre » à « Washnitzer-Monsky ».}

2.8. Il se pose la question, d'autre part, pour quels schémas abéliens les
énoncés 1º et 3º de 2.4, 2.5.\ et 2.6.\ sont valables, en déhors\note{sic.}
du cas déjà signalé : $S$ de caractéristique nulle. [J'aurais assez tendance à
croire que si $A$ est une VA définie sur un corps $k$ de car.\ $p>0$, alors ce
qu'on veut (savoir l'indépendance de $G(B)$ d'un relèvement infinitésimal
choisi $B$ de $A$) est vrai, pour tout anneau artinien de corps résiduel $k$,
\emph{si et seulement si} $A$ est « ordinaire » i.e.\ a sur la clôture
algébrique de $k$ le nombre maximum de points d'ordre $p$.\marginal{faux,
\uncertain{c'est jamais} vrai en car $p>0$ !}\note{au crayon, en biais dans la
marge de gauche, relié par un trait au $A$ de « si $A$ est une VA » ; le
crochet ouvrant avant « J'aurais » est à l'encre. Un long trait oblique au
crayon traverse la page de « un corps » jusqu'à la dernière ligne : il semble
annuler l'alinéa, ou le suivant.} Notes\note{sic.} qu'il n'est pas nécessaire
de supposer $k$ parfait (la question est manifestement géométrique), donc le
groupe \struck{formel} \add{$p$-divisible} de $A$ sur $k$ ne se décompose pas
nécessairement en produit d'un étale par un infinitésimal.\note{« $p$-divisible »
est tapé dans l'interligne au-dessus de « formel » biffé à la machine.} On
aurait alors, pour de telles VA au moins, la description très explicite des
relèvements infinitésimaux en termes du relèvement d'une filtration d'un module
à la Dieudonné. Cette conjecture semble raisonnable car on a fait ce qu'il
fallait pour éliminer des ennuis comme dans l'éxemple\note{sic.} de 2.7., d'autre
part cela s'harmonise fort bien avec le tapis Serre des relèvements canoniques
\add{(+ des modules)} de variétés abéliennes du type
« ordinaire ».\marginal{\ill{}}\note{« (+ des modules) » : « modules » tapé
dans l'interligne, « + des » et les parenthèses à l'encre. Une note au crayon,
de trois ou quatre mots, en biais dans la marge de gauche à hauteur de
« fallait pour éliminer », reliée par un trait au début de la ligne, n'est
pas lue ; un trait vertical au crayon longe les trois dernières lignes de
l'alinéa.}

Il serait évidemment tentant d'espérer qu'un énoncé du même type serait vrai
pour une base quelconque, pourvu que toutes les fibres de $A/S$ soient des VA
ordinaires. Mais cela est certainement faux avec cette généralité, même pour
des schémas elliptiques, essentiellement toujours pour les même raisons.\note{la
fin du mot, « sons. », est écrite à l'encre sous la ligne.}

\page{28}
\note{p.~27 de la lettre.}
Prenons en effet un schéma \add{$S$} modulaire à échelons \struck{$S$,} pour
variations de courbes elliptiques, et localisons en $\underline{Z}_p$, soit
encore $S$ le schéma lisse sur $\underline{Z}_p$ obtenu. On a sur $S$ le Module
loc libre de rang 2 $H^1_{\mathrm{DR}}(A/S)$, et quand on enlève de $S$ le
nombre fini de points fermés correspondants à Hasse nul, on aurait une
stratification \struck{\ill{}} rel.\ à $\underline{Z}_p$ sur ce Module. Pour
des raisons de pureté évidentes, celle-ci se prolongerait au dessus de $S$
tout entier, d'où une stratification sur le $H^1_{\mathrm{DR}}$ de la famille
modulaire de courbes elliptiques, en caractéristique $p$. Celle-ci serait alors
compatible avec frobenius (car elle l'est dans le complémentaire de l'ensemble
fini des points de Hasse nul), ce qui contredit nos réfléxions de 2.7.\note{le
« $S$ » après « schéma » est à l'encre, le « $S$, » tapé après « échelons »
est biffé à l'encre ; un mot tapé avant « rel. » est biffé à la machine. Un
long trait oblique au crayon traverse tout cet alinéa, qui prolonge celui de la
page précédente.}

Il n'est pas exclu entièrement, cependant, que en restant en car.\ $p>0$,
l'énoncé sur la non-variation \struck{\ill{}} infinitésimale de $G(A)$ avec
$A$, $A$ partout ordinaire, soit vrai. Cela impliquerait que sur l'ouvert des
valeurs « ordinaires » de l'invariant, le $H^1_{\mathrm{DR}}$ a une
stratification canonique, mais celle-ci ne se prolongerait pas à la courbe
modulaire toute entière. Mais je dois dire que ce drôle de comportement, où on
aurait une stratification naturelle en car.\ $p$ et une autre en car.~0, sans
qu'elles veuillent se recoller, semble assez canularesque.

[Les considérations précédentes font bien ressortir la nature « infinie » de
la notion de stratification, par contraste avec celle de connexion, malgré les
trompeuses apparences de la caractéristique nulle. Ainsi, sur le
$H^1_{\mathrm{DR}}$ de la famille modulaire sur $\underline{Z}$ \struck{pour}
\add{de} schémas elliptiques, il y a \struck{sur} \add{au dessus de} la fibre
générique $S_{\underline{Q}}$ de $S$ une stratification naturelle, mais nous
venons de voir (ou presque \ldots) que cette stratification ne s'étend pas en
une stratification sur $S_U$, $U$ un ouvert non vide de
$\operatorname{Spec}(\underline{Z})$. (Le « presque » provient du fait qu'il
n'est pas absolument clair si la stratification qu'on obtiendrait ainsi en
car.\ $p$ serait respectée par frobenius ; il faut absolument tirer au clair
cet exemple particulier !). C'est en un sens assez moral, puisque pour nous le
\add{module à} stratification doit jouer dans une large mesure le rôle d'un
faisceau $\ell$-adique, dont il partage\note{« de » et « au dessus de » sont à
l'encre, au-dessus de « pour » et « sur » biffés ; « module à » est tapé dans
l'interligne et appelé à l'encre (le « la » tapé est corrigé en « le ») ; le
crochet ouvrant est à l'encre. Le $\ell$ est tapé comme un X barré.}

\page{29}
\note{p.~28 de la lettre.}
également les propriétés de rigidité.]\note{crochet fermant à l'encre.}

Pour en revenir à l'alinéa précédent, concernant un schéma abélien
« ordinaire » en car.\ $p>0$, s'il n'est pas exclu que le $H^1_{\mathrm{DR}}$
admette une stratification canonique fonctorielle, il me semble cependant
exclu, malheureusement, que celle-ci provienne d'un cristal \struck{\ill{}} de
modules sur $S$, i.e.\ que ceci reste vrai en remplaçant $S$ par toute
extension infinitésimale, pas néc de car.\ $p$, \add{(tout au moins,} en
admettant que la connexion associée à la stratification en question soit la
connexion de Gauss-Manin\add{)}. En effet, appliquant une telle hypothèse au
schéma modulaire $S$ sur $\underline{Z}_p$ précédent, dont on enlèverait les
points de Hasse nul d'abord\add{, d'où $S'$}, on conclurait sauf erreur que
\struck{\ill{}} la stratification qu'on a en caractéristique nulle sur
$H^1_{\mathrm{DR}}$ se prolongerait à $S$ tout entier, car elle se
« recollerait » en un sens évident avec la \struck{\ill{}} stratification qu'on
aurait \struck{\ill{}} \add{au dessus du} complété $p$-adique de $S'$ (ce qui
doit impliquer le prolongement sur $S$ de façon assez formelle).\note{« (tout
au moins, » et la parenthèse fermante après « Gauss-Manin » sont à l'encre ;
« d'où $S'$ » et « au dessus du » sont tapés dans l'interligne et appelés à
l'encre ; les guillemets ouvrants de « recollerait » sont à l'encre. Les mots
biffés le sont à la machine.} Mais alors on aurait en car.\ $p$ une
stratification du $H^1_{\mathrm{DR}}$ au dessus du schéma modulaire tout
entier, Hasse nul inclus, ce qui est absurde comme on l'a déjà remarqué. Il
faut en conclure, hélas, que si les isocristaux, et le cas échéant les modules
stratifiés même en caractéristique \struck{nulle} $p>0$, ont des chances d'être
des outils convenables pour l'étude de familles de variétés abéliennes, celle
de cristal elle-même semble irrémédiablement trop fine, même en se restreignant
à des familles de VA « ordinaires » en car $p>0$. Elle peut tout au mieux se
prêter au cas d'un schéma de base réduit à un point, spectre d'un corps pas
nécessairement parfait, et on peut alors espérer les résultats les plus
satisfaisants en se bornant aux VA ordinaires. -- Pour que la notion de cristal
elle-même puisse être utilisée pour des schémas abéliens sur des bases plus
générales, il semble donc qu'il faille imposer aux familles envisagées des
restrictions très sérieuses, consistant à imposer la variation infinitésimale
du $G(A)$. Cela semble assez proche du point de vue de Serre, qui étudie les
variations de VA (éventuellement à multiplication complexe donnée) en imposant
à priori l'espace tangent (et l'action de la MC dessus) \ldots

\page{30}
\note{p.~29 de la lettre.}
\section*{\emph{Chapitre} 3\quad \emph{Remarques sur les groupes $p$-divisibles.}}

3.1. Bien sûr, en même temps que pour les VA, je dois en rabattre sur les
groupes formels. Je veux simplement signaler que la construction de $G(A)$ dans
2.6.\ doit pouvoir se paraphraser sans difficulté pour un groupe $p$-divisible
$\varnothing$ sur un préschéma $S$ à caractéristiques résiduelles égales au
même $p$.\note{le texte tape « p-disible », corrigé par un « v » à l'encre. Le
groupe $p$-divisible est noté par le signe tapé « Ø », rendu ici
$\varnothing$.} En effet les homomorphismes de ce groupe dans le groupe formel
associé à $\underline{G}_a$ sont triviales (en particulier, $\varnothing$ n'a
pas d'automorphismes infinitésimaux), d'autre part (sauf erreur) le faisceau en
\struck{groupes} \add{modules} des $H^2$ de $\varnothing$ à valeurs dans le
groupe formel associé à $\underline{G}_a$ (au sens du complexe du groupe
$\varnothing$, variante formelle), est un Module localement libre de rang
$g^*$, où $g^*$ est la dimension du groupe dual de $\varnothing$, soit
$\varnothing^*$ ; plus précisément, ce Module doit être canoniquement
isomorphe à $\underline{\mathrm{Lie}}(\varnothing^*)$ \add{$(=
\underline{t}_{\varnothing^*}.)$} \add{(}Tout ceci est suggéré par les
résultats de Tate et Lubin sur les modules de groupes formels ; je dois avouer
d'ailleurs que je n'ai pas essayé de tirer au clair ce qu'il faut entendre,
dans ce contexte, par « le groupe formel associé à $\underline{G}_a$ », et
s'il y a lieu de prendre le groupe formel purement infinitésimal.\add{)} Il
s'impose alors de paraphraser la construction de Serre, en regardant
l'extension \struck{\ill{}} universelle de $\varnothing$ par un groupe formel
vectoriel, qui sera ici le groupe formel (covariant) de
$\underline{t}_{\varnothing^*}$. Désignant par $G(\varnothing)$ cette
extension, son algèbre de Lie $H(\varnothing)$ sera une extension
\[
  (*)\qquad 0\to\underline{t}^{\vee}_{\varnothing^*}\to\underline{H}(\varnothing)
  \to\underline{t}_{\varnothing}\to 0\ .
\]
\note{« modules » est tapé au-dessus de « groupes » biffé à la machine ;
« $(=\underline{t}_{\varnothing^*}.)$ » est tapé dans l'interligne, avec le
signe $=$ et les parenthèses à l'encre, ainsi que les parenthèses qui
encadrent « Tout ceci \ldots\ infinitésimal. ». L'astérisque de
$\underline{t}_{\varnothing^*}$ dans « le groupe formel (covariant) de
$\underline{t}_{\varnothing^*}$ » et le $\vee$ de la suite exacte sont à
l'encre ; les flèches de la suite exacte sont repassées à l'encre.}

Bien entendu, $G(\varnothing)$ et par suite $\underline{H}(\varnothing)$
seront fonctoriels en $\varnothing$, et compatibles avec changement de base.
Alors qu'il est certainement vrai, dans un sens qu'il conviendrait de préciser,
que \struck{\ill{}} $G(\varnothing)$ varie « moins que $\varnothing$ », quand on
fait varier $\varnothing$ infinitésimalement dans une famille, il n'est pas
vrai pour autant que $G(\varnothing)$ soit indépendant de telles variations (à
l'exception, probablement, de celles du premier ordre), i.e.\ puisse être
considéré comme provenant d'un cristal en groupes \struck{$p$-divisibles}
\add{$\pm$ formels}. Cela sera le cas, seulement, sans doute, quand
$\varnothing$ sera extension d'un groupe \struck{formel} ind-étale par un
groupe $p$-divisible torique\add{, et $S$ réduit à un point (?).}
\add{(}Dans le cas général, il semble intéressant d'étudier les
variations\note{« $\pm$ formels. » est à l'encre au-dessus de « p-divisibles. »
biffé à l'encre ; « formel » est biffé à la machine ; « , et $S$ réduit à un
point (?). » et la parenthèse ouvrante qui suit sont à l'encre, dans
l'interligne.}

\page{31}
\note{p.~30 de la lettre.}
infinitésimales de $\varnothing$, quand on se fixe celles de $G(\varnothing)$ à
l'aide d'un cristal en groupes $\underline{G}$.

3.2. En tous cas, l'extension $(*)$ semble un invariant intéressant du groupe
$p$-divisible envisagé. Il subsiste, par passage à la limite, en passant
\struck{à} au cas où $S$ est \struck{un schéma formel noethérien} par exemple
le spectre d'un anneau $A$ noethérien $j$-adique\note{sic, pour
$J$-adique.} séparé et complet, où $J$ est un idéal tel que $A/J$ soit à
caractéristiques résiduelles égales à $p$. On trouve par exemple un bon
invariant quand $A$ est un anneau de valuation discrète complet, éventuellement
d'inégales caractéristiques, à caractéristique résiduelle $p$. Dans le cas où
le groupe formel\add{,} réduit a la structure triviale \struck{\ill{}}
mentionnée plus haut dans 3.1., et que le corps résiduel $k$ est parfait, d'où
un module de Dieudonné $M$ sur $W$ (de rang $g+g'$), il semble donc, d'après
la remarque faite dans 3.1., que la donnée d'un relèvement du groupe
$p$-divisible qu'on a sur $k$ en un groupe $p$-divisible sur $A$, soit
entièrement équivalente à la donnée d'un relèvement de la filtration de
$M\otimes_W k$ donnée par $(*)$, en une filtration de type $(g,g^*)$ de
$M\otimes_W A$, \struck{Comme} Il serait bien intéressant d'essayer de
déterminer, en termes d'une telle donnée, quel est le module de Tate
correspondant à la fibre générique de $\varnothing$ ! On aimerait préciser
également, pour des groupes $p$-divisibles plus généraux, quelle quantité
d'information est liée à l'extension $(*)$\add{,} qui remplace ici le $H^1$ de
De Rham.\note{les biffures sont à la machine ; les virgules après « formel »
et après « $(*)$ » sont à l'encre. Le rang est écrit $g+g'$ ici et le type
$(g,g^*)$ deux lignes plus bas : tel quel. La phrase « Dans le cas où le
groupe formel, réduit \ldots » est laissée telle qu'elle est tapée.}

\note{le tapuscrit s'arrête ici, au milieu de la p.~30 de la lettre, sans
formule finale ni signature ; le reste de la page est blanc. Ce qui suit dans
le dossier est d'une autre nature : des notes manuscrites de sa main.}

\section*{Réduction des courbes elliptiques et $T_p(E)$}
\note{titre écrit au crayon, de sa main, en haut à droite d'une feuille par
ailleurs blanche (p.~32 des archivistes), qui sert de chemise aux notes
manuscrites qui suivent ; il ajoute : « (cf.\ aussi livre de Serre "Abelian
$l$-adic representations") ». Ces notes ne font plus partie de la lettre.}

\page{33}
Courbes elliptiques sur un \struck{anneau de} \uncertain{corps} local à
car.\ résiduelle $p>0$ : étude infinitésimale.

$A$ anneau de val.\ discrète complet, corps des fractions $K$, corps résiduel
$k$ de car $p>0$.

$E_K$ courbe elliptique sur $K$, \struck{\ill{}}.\note{trois ou quatre mots
biffés, dont peut-être « module » à la fin.}

1º) \emph{$K$ de car 0} i.e.\ $T_p(E_{\overline{K}})$ est \emph{de rang 2}
\marginal{[\ill{} et $T_p(E_{\bar{k}})$ \uncertain{ét} $=T_p(E_{\bar{k}})$]}\note{au
crayon, à droite, entre crochets ; lecture très douteuse.}
On veut étudier les structures possibles pour ce module galoisien.

\quad a) $E_K$ a « bonne réduction » i.e.\ $E$ provient d'une courbe elliptique
$E$ sur $A$, \emph{et} $E_k$ est « ordinaire ».

Alors $T_p(E_K)$ est \add{can.} extension d'un \struck{\ill{}} $L$ de rang~1,
par un \struck{autre} \add{\ill{}} \ill{} que le \uncertain{foncteur}
\ldots\note{« can. » (canoniquement ?) est écrit au-dessus de la ligne et appelé
avant « extension » ; au-dessus de « par un », une insertion de trois ou
quatre mots, soulignée, n'est pas lue ; la suite de la ligne n'est pas lue
davantage.}
\struck{Hom}
\[
  \check{L}\otimes T_p(\mathbf{G}_m)=\check{L}[1]\ .
\]
\uncertain{De plus},
\[
  0\to \check{L}(1)\to T_p(E_K)\to L\to 0\ ,
\]
l'extension\struck{, qui} est donnée par un élément de
\[
  H^1\bigl(K,\ L^{-2}\otimes_{\mathbf{Z}_p}T_p(\mathbf{G}_m)\bigr)\ ,
\]
qui dans le cas où $A$ est à corps résiduel sép.\ clos (dans le
\uncertain{cas} $L$ \uncertain{trivial} sur $A$) \uncertain{s'identifie à}
\[
  H^1(K,T_p(\mathbf{G}_m))\otimes L_0^{-2}\simeq
  \bigl(\varprojlim_{p^n} K^*/K^{*p^n}\bigr)\otimes_{\mathbf{Z}_p} L_0^{-2}
  \simeq \mathrm{Hom}\bigl(L_0^{2},\widehat{K^*}\bigr)
\]
\note{dans le $H^1$, deux ou trois signes sont hachurés avant $L^{-2}$ ; un
$*$ est barré au-dessus du $L$ de la suite exacte. Le quotient
$K^*/K^{*p^n}$ est restitué : on lit « $\varprojlim K^*$ », un signe noirci,
et $p^n$ en indice ; « $K^*$ » est aussi écrit, puis biffé, au-dessus de la
ligne. Il note $L^{-2}$ pour $\check{L}^{\otimes 2}$. Sous la formule :
« $\mathbf{Z}_p$, $L_0$ la valeur en $k$ » (lecture douteuse), qui
définit sans doute $L_0$.}

\page{34}
Si \uncertain{cet} \struck{extension provient \ill{}} élément est \ill{}
« \ill{} » i.e.\ est dans le sous-groupe
\[
  \mathrm{Hom}_{\mathbf{Z}_p}(L^2,\underbrace{A^{**}}_{\text{Einseinheiten}})
  \simeq L^{-2}\otimes_{\mathbf{Z}_p}A^{**}\ .
\]
\note{« Einseinheiten » (unités principales), en allemand, sous $A^{**}$, avec
une accolade.}

[\emph{N.B.} Le fait \ill{} \ill{} \ill{} Einseinheiten \uncertain{est} \ill{}
\uncertain{détermine} \ill{} mod \struck{\ill{}} \add{\uncertain{exponentielles}}
\ill{} \ill{} \ill{} \ldots\ $p$-adiques \ldots]\note{passage entre crochets,
souligné en partie ; le mot au-dessus du mot biffé est inséré par un
signe~$\vee$.}\marginal{\emph{N.B.} L'extension $T_p(E_K)$ \uncertain{peut
être triviale} (\ill{} de Serre) ou \ill{} : \ill{} \ill{} \ill{} \ill{}
d'ordre \uncertain{fini} \ldots\ « \ill{} »}\note{note de sa main, à droite
de l'alinéa précédent, en quatre lignes serrées.}\marginal{En tous cas,
\ill{} \uncertain{opération} \ill{} \ill{} [\ill{}]}\note{dans la marge de
gauche, en biais.}

\quad b) $E_K$ a bonne réduction, et $E_k$ \ill{} i.e.\ est non ordinaire.
J'ai envie de conjecturer que \uncertain{dans} la représentation de $\pi_1(K)$
dans $T_p(E_{\overline{K}})$ est $\frac12$-simple, \uncertain{voire} que
l'image de $\pi_1(K)$ dans $\mathrm{Aut}(T_p(E_{\overline{K}}))$ est un
\uncertain{sous-groupe} \ill{}ment \ldots\ ??\note{un grand crochet dans la
marge de gauche embrasse cette conjecture. « $\frac12$-simple » :
semi-simple, dans sa notation.}

(a$'$+b$'$) $E_K$ n'a pas \add{néc} bonne réduction, mais l'invariant de $E$
est entier, i.e.\ après extension finie de $K$, on obtient bonne réduction.
Ceci \uncertain{explique} \ill{} \ill{} , a$'$) réduction \ill{} (la base
\ill{}, i.e.\ \ill{} \ill{} ordinaire après ext.\ de $K$). Alors
l'\uncertain{application} \ill{} \ill{} \ill{} \ill{} \ill{}.\add{N.B.} de
$\pi'(K)$ est « \ill{} ». b$'$) \ill{} \ill{} \ill{} \ill{}. Alors
\uncertain{si} la conjecture \uncertain{dans} b) est \ill{}, la conclusion de
cette conjecture \uncertain{vaut} aussi dans le cas b$'$).\note{les deux
dernières lignes, au pied de la page, sont d'une autre encre, brune.}

\page{35}
que cet élément a une image \emph{non nulle} dans $\mathbf{Z}_p$, donc que
l'extension ne provient pas d'une extension de groupes $p$-divisibles sur
\emph{$A$}.

\emph{Question} Quel est l'élément obtenu de $\mathbf{Z}_p$ ! est-ce toujours
un entier, resp.\ un élément $>0$ ?

2º) \emph{Cas} \uncertain{car} $K\neq 0$ \emph{donc} car $K=p$.

\struck{\ill{}} Il y a lieu d'\uncertain{étudier} \ill{} $T_p(K_s)$, qui ne
s'explicite \uncertain{plus} en \ill{} en termes de \uncertain{modules}
exclusivement \ill{} galoisiens.\note{« $T_p(K_s)$ » : lecture douteuse ;
on attend $T_p(E_{K_s})$.}

\quad A) Cas général (bonne réduction). a/ Si \uncertain{quand} \ill{}
$E_K$ est ordinaire, donc \ill{} ordinaire, alors $T_p(E_K)$ est une
extension d'un $L_K$ par $\check{L}_K\otimes T_p(\mathbf{G}_m)$, et
l'extension est donnée \add{(qui \ill{} \ill{} \ill{})} \uncertain{dans}
$L^{-2}\otimes A^{**}$ (par \ill{} \ill{} invariant \ill{} dans
$L^{-2}\otimes$\note{la page s'arrête sur ce $\otimes$ ; un trait fléché
horizontal précède « $T_p(E_K)$ » en tête de ligne.}

\page{36}
\emph{NB} Si l'invariant $j$ est entier, \uncertain{alors} on peut
conjecturer que \struck{la réduction de} $E_K$ a bonne réduction ssi il en est
de \uncertain{même} de $T_p(E_K)$, le groupe $p$-divisible
associé.\marginal{?}\note{point d'interrogation dans la marge de gauche ; un
signe $=$ est biffé devant « $T_p(E_K)$ ».}

\quad c) L'invariant $j$ n'est \add{\uncertain{pas}} entier. Alors le
\emph{\uncertain{modèle} de Néron} \add{$E$} est tel que $E^0_k$ est
\add{\uncertain{forme de}} $\mathbf{G}_m$. Par suite, on conclut que dans
$p^{\nu}E$ \add{[quitte à faire une ext.\ finie sur $A$ au
\uncertain{plus}]} il y a un sous-groupe canonique \struck{\ill{}} forme de
$\mu_{p^{\nu}}$. On trouve ainsi dans $T_p(E_K)$ un sous-gr., qui est une
forme \struck{finie de} \add{\uncertain{tordue}} de $T_p(\mathbf{G}_m)$
\add{[donc aussi \ill{} $T_p(\mathbf{G}_m)$ \ill{} \ill{} pas]}. Donc si $k$
est sép.\ clos, on trouve que $T_p(E_K)$ admet un sous-module isomorphe [de
façon canonique \uncertain{au} signe \uncertain{près}] à $T_p(\mathbf{G}_m)$
\add{$=\mathbf{Z}_p(1)$}, \uncertain{le} quotient étant
\uncertain{canoniquement} isomorphe à $\mathbf{Z}_p$ [car
$\det T_p(E_K)\simeq\mathbf{Z}_p(1)$]. L'extension est ici déterminée par un
élément canonique de $\widehat{K^*}$, ext.\ de $\mathbf{Z}_p$
[\uncertain{provenant} de $K^*/A^*$] par $A^{**}$. Je conjecture\note{les
ajouts sont écrits au-dessus de la ligne et appelés par un trait ; un mot
souligné deux fois, « \uncertain{tordue} », est écrit au-dessus de « finie
de » biffé. La page s'arrête sur « Je conjecture » ; la suite manque.}

\section*{Gribouillis de Tate sur sa théorie}
\note{titre écrit de sa main, en haut à droite d'une feuille par ailleurs
blanche (p.~37 des archivistes), qui sert de chemise à la feuille suivante.}

\page{38}
\note{feuille d'une autre main, à l'encre très pâlie, en anglais : d'après
la chemise, notes de Tate sur sa théorie (hypothèse « $L/K$ Galois », anneau
des entiers $R_L$, groupe $\mathcal{G}=\mathrm{Gal}(L/K)$, un « Thm. » reliant
$\widehat{L}^{\mathcal{G}}=\widehat{K}$ à la nullité de
$H^1(\mathcal{G},R_L)$, suites exactes en $H^1(\mathcal{G},R_L)$). Elle n'est
pas de lui et n'est pas reproduite ; on n'y voit aucune intervention de sa
main.}

\section*{Questions diverses en théorie de Tate-Hodge}
\note{titre écrit de sa main, en haut à droite d'une feuille par ailleurs
blanche (p.~39 des archivistes). Le crayon des archivistes porte 50 sur cette
chemise et 51 sur la feuille suivante, alors que le fac-similé les place en
39 et 40 : la numérotation suivie ici est celle du fac-similé.}

\page{40}
\note{écrit à l'encre au verso d'une page dactylographiée étrangère au
dossier (texte tapé en transparence, non reproduit).}
\emph{Questions} a) \uncertain{Théorème} \ill{} \ill{} (fonctoriel --
multiplicatif en $X$, et commutant aux transport de structure)
\[
  H_{\mathrm{DR}}(X_{\mathbf{C}})\simeq H_{\mathrm{Hdg}}(X_{\mathbf{C}})
\]
compatible avec filtrations \add{\uncertain{du moins}} une fois donné un
\add{générateur} $\varepsilon\in T_\ell(\mathbf{C}^*)$ ?\note{« générateur »
est écrit au-dessus de la ligne et appelé par un trait ; les deux mots au-dessus
de « filtrations » sont de lecture douteuse.}

\quad b) Tout $\mathbf{Q}_\ell(\pi)$-module simple, de dim finie sur
$\mathbf{Q}_\ell$, est-il contenu dans un $\mathbf{C}(i)$ ?
\uncertain{Cas} \uncertain{où} il \ill{} \ill{} fini \ill{} $M$ et un $M(i)$,
pour $i\neq 0$ ??\note{« $\mathbf{C}$ » est écrit sur un autre signe ;
lecture douteuse de toute la phrase.}

\quad c) Y-a-t-il d'autres conditions sur les
\[
  M=H^i(X_{\mathbf{C}},\mathbf{Z}_\ell)\otimes_{\mathbf{Z}_\ell}\mathbf{Q}_\ell
\]
que celles résultant de la conjecture de Tate ?

\quad d) \emph{Construction} des \uncertain{morphismes} de Tate
\[
  H^{pq}_{\mathrm{D}}(X)_K\simeq H^q(X,\Omega^p_X)
\]
\note{l'exposant du premier membre est de lecture douteuse.}

\quad e) Lien de ces structures avec les connexions canoniques ??

\quad f) Lien avec Washnitzer-Monsky ?

\end{document}
